% TARGET STATEMENT ONLY: night:P1.19
\begin{theorem}\label{night:P1.19}
\StatementRecord{night:P1.19}{proved}{}{\cref{night:P1.11}, \cref{night:P1.12}, \cref{night:P1.9}; \cite[Theorem~3.43, p.~254]{Hatcher2002}}
There is an independent computation from the full global cellular
complex.  Its chain ranks in degrees zero through eight are
\[
 (4,26,87,162,200,156,79,23,3),
\]
and its differentials, after the explicit integral chain equivalences
below, are the matrices in
\path{code/current_audits/audit_corrections/code/reconstruct_V_global_06.json}.
The accompanying program reconstructs them from the displayed marked
maps; the separate saved-output verifier supplies their integral
cycle and incoming-map checks and the basis comparisons below.
Their integral Smith calculation gives
\[
 H_k(M;\mathbb Z)=
 \begin{cases}
 \mathbb Z,&k=0,8,\\
 \mathbb Z^4,&k=4,\\
 0,&\text{otherwise}.
 \end{cases}
\]
This calculation uses local Poincar\'e--Lefschetz duality of the
smooth cusp cap to determine its secondary attaching class.  It does
not use global Poincar\'e duality or the homology conclusion of P1.12.
In particular, the map obtained by setting the whole cusp suspension
chain term to zero is not the actual relative cusp model; its precise
missing term is accounted for here.
\end{theorem}

% CITED INPUT, WITH ITS PROOF: night:P1.10
\begin{theorem}\label{night:P1.10}
\StatementRecord{night:P1.10}{proved}{}{\cref{night:P1.11}, \cref{night:P1.9}}
For every real parameter $c>-h_{\min}$ the corrected construction gives
a compact smooth complex fourfold $M_c$ and a proper holomorphic map
to $\mathbb P^1$ with the two specified multiple fibres and the
index-one Kuhn cusp.  With its complex orientation,
$\pi_1(M_c)=1$ and $\chi(M_c)=6$.
\end{theorem}
\begin{proof}
\begin{enumerate}[label=\arabic*.]
\item Evaluation of $\Lambda$ on $F_c$ maps $\Lambda_{\mathbb R}$
isomorphically onto the underlying real bundle of $F_c^\vee$:
the kernel is zero exactly by $F_c\cap\bar F_c=0$.
The image is consequently a full discrete lattice.  Quotient by its
holomorphic period sections gives a holomorphic torus bundle on the
nonspecial base.

\item On a finite root chart $u^{m_i}=t_i$, divide the pullback torus bundle by
\[
(x,u)\longmapsto(M_ix+v_i/m_i,\zeta_{m_i}u).
\]
Freeness proves smoothness, including the central fibre.  The special
divisor is $m_i$ times the reduced free quotient.  On the punctured
root disk translation by
$-(\log u/(2\pi i))v_i$, interpreted by its period section, identifies
this affine quotient with the original Jacobian bundle.  A change of
logarithm gives an integral lattice period; replacing $u$ by
$\zeta_{m_i}u$ cancels exactly the affine term $v_i/m_i$.
This proves the holomorphic gluing and its sign.
At the cusp use P1.9 and the exponential period map.  The zero section
$\zeta=0$ extends to $z=1$ in the dense torus of the component marked
by the vertex zero.  The resulting map is proper by a finite compact
cover of the base by the three local pieces and a compact remainder
with torus fibres.  All total-space charts and their overlaps have
already been specified, so this produces a smooth compact complex
fourfold, and in particular an almost complex structure.

\item For completeness we record every van Kampen relation.  The main
bundle has its zero section and fundamental group
$\Lambda\rtimes\langle g_1,g_2\rangle$.
The finite pieces impose
$g_i^{m_i}=v_i$ and $g_ix^ug_i^{-1}=x^{M_iu}$; the logarithmic gluing
above supplies these signs.  The cusp zero section kills
$g_0=(g_1g_2)^{-1}$ and the cusp collapse kills $\operatorname{im}N$.
The latter assertion follows directly from the toric moment model in
P1.11: angular cycles $a_j$ collapse and the base cycles $b_j$ survive.
Its fundamental group is $\mathbb Z^3$, since the simply connected
complete toric components and their connected toric intersections have
dual complex $\mathbb R^3/\mathbb Z^3$; alternatively, van Kampen on
the moment cells kills the angular loops generated by the primitive
facet normals and retains exactly the three base loops.
Thus the full presentation is
\[
\left\langle \Lambda,g_1,g_2\ \middle|\
 [\Lambda,\Lambda]=1,\quad
 g_ix^ug_i^{-1}=x^{M_iu},\quad
 g_1^5=v_1,\ g_2^4=v_2,\ g_1g_2=1,\ \operatorname{im}N=0
\right\rangle.
\]
The last relation is redundant after the preceding ones but records
the actual cusp map.  Since $g_2=g_1^{-1}$, the two induced actions are
inverse; their orders divide both five and four and are therefore one.
The remaining fibre group is $\mathbb Z$ through $\ell$.  With its
generator $z$, the additive relations are
$5g_1-z=0$, $4g_2+z=0$, $g_1+g_2=0$.
Their matrix
\[
\begin{pmatrix}5&0&-1\\0&4&1\\1&1&0\end{pmatrix}
\]
has determinant $-1$, proving triviality.
Euler additivity gives $\chi(M_c)=6$: smooth tori and finite free
torus quotients have Euler number zero, and P1.9 gives the remaining
contribution.

\end{enumerate}
\end{proof}

% CITED INPUT, WITH ITS PROOF: night:P1.11
\begin{theorem}\label{night:P1.11}
\StatementRecord{night:P1.11}{proved}{}{\cref{night:P1.9}}
The index-one central cusp fibre $C$ has torsion-free integral homology
with ranks $(1,3,9,9,10,3,1)$.
The collapse $\kappa:T^6\to C$ has primitive image in every degree;
it is onto except in degree four, where its cokernel is $\mathbb Z$.
Its image ranks are $(1,3,9,9,9,3,1)$.
In every degree its kernel is the saturation of
$(\bigwedge^kM_0-I)\bigwedge^k\Lambda$.
\end{theorem}
\begin{proof}
\begin{enumerate}[label=\arabic*.]
\item We first remove the unit matrix $C(q)$ topologically, by an explicit
deformation, without changing the integral marking.  Use the family
$C_u(q)=(1-u)C(q)+u iI_3$, followed by the reverse of the same
family starting at $C(q)=0$.  The two flows compose to a marked
comparison with $C(q)=0$.  The action and properness estimates
of P1.9 are uniform on both parameter intervals.  Local quotient charts form a smooth
family of pairs, with the local equation $q=x_0x_1x_2x_3$ independent
of $u$.  Local lifts of $\partial/\partial u$, tangent to every
coordinate stratum and annihilating $q$, can be patched by a partition
of unity.  Over a smaller compact disk their flow exists for the whole
interval.  This gives a stratified diffeomorphism of the central
fibres and a diffeomorphism of the nearby fibres, preserving the
period marking.  The same construction on the real oriented blowup
preserves the collapse map.  We may therefore calculate at $C(q)=0$.

\item Here is a concrete topological model.  Put $G=4I_3-\mathbf1\mathbf1^t$.
Its positive eigenvalues are $1,4,4$.  The Voronoi polytope in dual
coordinates is cut out by
\[
|v^ty|\leq \tfrac12v^tGv
 \quad (0\ne v\in\{0,1\}^3).
\]
The value on a binary vector with $k$ nonzero entries is $k(4-k)$.
These fourteen primitive facet normals give the star of a vertex in
the Kuhn triangulation.  One checks its Delaunay tetrahedra by the
successive chains
$0,e_i,e_i+e_j,e_1+e_2+e_3$ for the six permutations $(i,j,k)$;
solving their three perpendicular-bisector equations and the fourteen
inequalities verifies precisely their common faces.  The translates
of this Voronoi polytope by $G\mathbb Z^3$ tile $\mathbb R^3$.

\item For a simplex $\sigma$ with vertices $v_0,\ldots,v_r$, put
$L_\sigma=\sum_i\mathbb Z(v_i-v_0)$ and
$D_\sigma=\bigcap_{v\in\sigma}P_v$, where $P_v=P+Gv$.
The component torus has lattice $A=\mathbb Z^3$; the common-stratum
marking is the explicit map
\[
 (A\oplus\mathbb Z)/\sum_i\mathbb Z(v_i,1)
 \longrightarrow A/L_\sigma,\qquad
 [(a,k)]\longmapsto a-kv_0.
\]
The Kuhn unimodularity makes $L_\sigma$ primitive.

We choose compatible moment maps as follows.  The polytope $2P_v$
is integral, since its vertex facet normals are unimodular and its
facet constants are integral.  Embed the component $X_v$ using the
characters $\chi^{-m}$ for
$m\in\mathcal A_v=2P_v\cap A^\vee$, and define
\[
 \mu_v([Z_m])=
 \frac{\sum_{m\in\mathcal A_v}m|Z_m|^2}
 {2\sum_{m\in\mathcal A_v}|Z_m|^2}.
\]
The sign is specified: the ray $w-v$ tends to the face on which
$\langle w-v,-m\rangle$ is minimal, hence to the upper common
Voronoi face of $P_v$ and $P_w$.  Thus the character polytope is
$-2P_v$, and $\mu_v$ is minus one-half its usual moment map.
The monomial moment theorem
\cite[equations (12.2.3)--(12.2.4) and Theorem 12.2.5, p.~571]{CoxLittleSchenck2011}
identifies the nonnegative toric part with $P_v$.
For the squared weights here, the intervening map on that part is
$[a_m]\mapsto[a_m^2]$, with inverse $[b_m]\mapsto[\sqrt{b_m}]$:
both preserve the homogeneous binomial relations and positive common
projective rescaling.  Write $s_v:P_v\to X_v$ for the resulting
continuous positive section.

On a shared face put
$\mathcal A_\sigma=2D_\sigma\cap A^\vee$ and choose
$m_0\in\mathcal A_\sigma$.
For $m\in\mathcal A_\sigma$, the difference $m-m_0$ annihilates
$L_\sigma$.  The ratio $Z_m/Z_{m_0}$ is consequently the restriction
of the same character
\[
 \bigl(-(m-m_0),\langle m-m_0,v\rangle\bigr)
       \in(A\oplus\mathbb Z)^\vee
\]
for every vertex $v$ of $\sigma$.  At $C=0$ all deck unit multipliers are one, and the toric gluing
identifies these character functions exactly.  On either component
the restricted moment map is therefore the identical expression
\[
 \frac{\sum_{m\in\mathcal A_\sigma}
       m|\chi^{-(m-m_0)}|^2}
 {2\sum_{m\in\mathcal A_\sigma}|\chi^{-(m-m_0)}|^2}.
\]
The equality persists on every smaller face by restricting the index
set.  Positive character values also agree, so the sections $s_v$
glue on the entire shared orbit closure.  The angular maps are the
displayed quotient homomorphisms $A\to A/L_\sigma$.

The period deck map is $(a,k)\mapsto(a+k\nu,k)$, sending $v$ to
$v+\nu$.  Its component marking is unchanged because
$a+k\nu-k(v+\nu)=a-kv$.
Moreover $\mathcal A_{v+\nu}=\mathcal A_v+2G\nu$; the common
character factor cancels in projective ratios.  Hence
\[
 \mu_{v+\nu}(T_\nu p)=\mu_v(p)+G\nu,
 \qquad s_{v+\nu}(y+G\nu)=T_\nu s_v(y).
\]
These give an explicit well-defined map
\[
 \begin{gathered}
 \mathcal Q=
 \bigl((\mathbb R^3/G\mathbb Z^3)\times
       (\mathbb R^3/\mathbb Z^3)\bigr)/\!\sim
 \longrightarrow C,\\
 [y,\theta]\longmapsto e^{2\pi i\theta}\,s_v(y)
 \quad(y\in P_v).
 \end{gathered}
\]
where over $\operatorname{relint}D_\sigma$ the relation collapses
the subtorus $L_\sigma\otimes(\mathbb R/\mathbb Z)$.
On each stratum the inverse takes its moment value and its angular
character class in $T_A/T_{L_\sigma}$; the latter is obtained by
$z\mapsto z/|z|$ on its nonzero characters.
The map is a continuous bijection, and its source is compact and
$C$ Hausdorff, so it is a homeomorphism.
The angular identity $[y,0]$ gives the continuous section, with
the same face and deck compatibilities just proved.

\item The actual nearby-fibre collapse is the quotient map in this model.
Give the semistable family its divisorial logarithmic structure.
The local equations $q=x_0x_1x_2x_3$ and their face restrictions
make it fine, saturated, log smooth and vertical over the standard
logarithmic disc.  The family is proper by P1.9.
In its logarithmic central fibre, retain all arguments of the $x_j$
with their sum fixed by $\arg q$.  The integral angular lattice over
a face with vertices $v_0,\ldots,v_r$ is the kernel of the height
map on $\mathbb Z^3\oplus\mathbb Z$.  Forgetting the arguments of
the zero coordinates induces exactly
\[
 \mathbb Z^3\longrightarrow
 (\mathbb Z^3\oplus\mathbb Z)/
       \sum_i\mathbb Z(v_i,1)
 \ \cong\ \mathbb Z^3/L_\sigma.
\]
The last isomorphism sends the class of $(a,k)$ to
$a-kv_0\pmod {L_\sigma}$; its inverse sends $a$ to $(a,0)$.
Thus no finite index enters this angular quotient.
The polar maps agree under the specified integral monomial overlaps.
At $C(q)=0$ the deck multipliers have positive real values, so the
logarithmic fibre at angle one is the product of the base and angular
three-tori in item~3, with their given markings.

The mapping-cylinder comparison for a proper vertical log-smooth
family identifies this logarithmic forgetful map with the actual
nearby-fibre collapse, compatibly with the central strata
\cite[formula (3-3-1), p.~478; Theorem 4.1.1 and Lemmas
4.1.2--4.1.3, pp.~480--482]{ffAchingerOgus20}.
Its angular cycles are $a_j$ and its base cycles are $b_j$.
In these column-basis conventions monodromy fixes each $a_j$
and sends $b_j$ to $b_j+a_j$, as specified by $Nb_j=a_j$.
The stratified deformation in item~1 transports this actual map
back to the original bounded period matrix.

\item Filter $\mathcal Q$ by the dimension $d$ of its dual base cells.
Over the interior of the dual cell of a $p$-simplex, $d=3-p$ and the
angular quotient has lattice
$Q_\sigma=\mathbb Z^3/L_\sigma$ of rank $d$.
Excision for the cell and its boundary gives the relative chain
groups of the resulting homology spectral sequence:
\[
E^1_{d,q}=\bigoplus_{\dim\sigma=3-d}\bigwedge^qQ_\sigma.
\]
Orient dual cells by the orientation of the base and the ordered
simplex.  The differential is the signed quotient map for each
codimension-one coface.  More explicitly, with integral projection
matrices $Q_\sigma$ and chosen right inverses $S_\sigma$,
\[
D_{p,q}[\tau,\sigma]=
 \sum_{\sigma=\partial_j\tau\ {\rm mod}\ \mathbb Z^3}
 (-1)^j\bigwedge^q(Q_\tau S_\sigma).
\tag{1.11.1}
\]
All $L_\sigma$ are primitive because the tetrahedra are unimodular.
Thus these projections and right inverses are integral.  Formula
(1.11.1) specifies every integer entry, including multiple incidences
under translation identifications.  The exact reconstruction in
\path{code/current_audits/audit_corrections/code/reconstruct_V_cusp_06.py} enumerates all $26$ simplex classes and
$74$ codimension-one incidences. Its JSON output supplies every integral
quotient and right inverse, each differential, cycle basis, incoming
matrix, homology projection and marked collapse matrix. All Smith
identities and unimodular basis changes are checked by exact multiplication.

\item The nonzero $E^2_{d,q}$ ranks are
\[
\begin{array}{c|rrrr}
q\backslash d&0&1&2&3\\\hline
0&1&3&3&1\\
1&0&6&8&3\\
2&0&0&7&3\\
3&0&0&0&1
\end{array}
\]
and every group is free. The reconstructed incoming Smith matrices
have only unit nonzero factors and reproduce every entry of this table.
Support leaves only the possible higher differential
$d_2:E^2_{3,0}\to E^2_{1,1}$.
The angular identity section splits the base projection as a map of
filtered spaces, so the base fundamental class survives.  This kills
that differential.  All subsequent differentials have zero source or
target.  The filtration extensions are extensions of free abelian
groups and therefore split, proving the asserted integral homology.

\item The product-to-quotient map induces on $E^1$ the exterior powers of
the quotient projections.  The program \path{code/current_audits/audit_corrections/code/reconstruct_V_cusp_06.py} constructs the induced maps
on integral cycle and quotient bases from the face projections and the
Poincar\'e-dual Alexander--Whitney formula in step 9. Its complete
integral output is \path{code/current_audits/audit_corrections/code/reconstruct_V_cusp_06.json}.
Their images are primitive, and the only nonzero graded cokernel
is $\mathbb Z$ at $(d,q)=(2,2)$.  The summed image ranks are
\[
             r_k=(1,3,9,9,9,3,1).
\]
The same exact exterior-power calculation gives these ranks for
$\ker(\bigwedge^kM_0-I)$.
The actual collapse satisfies
$\kappa_*(\bigwedge^kM_0-I)=0$, because a monodromy circuit
has the same retraction to the central fibre in the mapping-cylinder
comparison.  Therefore its rank is at most $r_k$.
A filtered map has rank at least the sum of the ranks of its
graded maps; hence the actual rank is exactly $r_k$.

Choose integral bases compatible with the source and target
filtrations; their graded groups are free.  In each graded block
choose the unit Smith pivots of the computed map.  The corresponding
$r_k$ rows and columns of the actual filtered map form a block
triangular submatrix with unimodular diagonal pivot blocks.
Its determinant is $\pm1$.  Since the full map has rank $r_k$,
these unit pivots prove that its actual image is primitive.
Thus the actual map is onto in every degree except four, where its
cokernel is the free group of rank $10-9=1$.
This rank argument also rules out additional lower-filtration images
of source classes in the kernels of the graded maps.

\item For a fully marked certificate we represent a base cycle $b_J$ by
Poincar\'e duality from the Alexander--Whitney cup cocycle on the
complementary coordinates.  On an ordered $p$-simplex its coefficient
is
\[
\operatorname{sgn}(J,J^c)
 \prod_{r=1}^p(v_r-v_{r-1})_{(J^c)_r}.
\]
Tensoring with an angular exterior vector and applying (1.11.1)
gives the marked matrices exported in
\path{code/current_audits/audit_corrections/code/reconstruct_V_cusp_06.json}. Their source basis is the ordered exterior
basis $(a_1,a_2,a_3,b_1,b_2,b_3)$; the output gives the integral
target homology projections and all Smith basis maps.
They annihilate $\bigwedge^kM_0-I$ exactly.  The actual collapse also
annihilates that operator, since turning the smoothing parameter
does not change its image in the central fibre.  The computed ranks
equal the ranks of the monodromy coinvariants; primitivity therefore
identifies both kernels with the saturation of its image.
Consequently the exported matrices are matrices of the actual
collapse in suitable integral target bases.  This last argument also
eliminates any ambiguity from choosing splittings of the filtration.

\end{enumerate}
\end{proof}

% CITED INPUT, WITH ITS PROOF: night:P1.12
\begin{theorem}\label{night:P1.12}
\StatementRecord{night:P1.12}{proved}{}{\cref{night:P1.1}, \cref{night:P1.9}, \cref{night:P1.10}, \cref{night:P1.11}; the actual marked evaluation and finite logarithmic collar maps; \cite[Theorem 4.1.1 and Lemmas 4.1.2--4.1.3, pp.~480--482]{ffAchingerOgus20}; \cite[Theorem 3.2, p.~195; Theorem 3.30, p.~241; Theorem 4.32, p.~366]{Hatcher2002}}
The fourfold $M_c$ has torsion-free integral homology
\[
H_k(M_c;\mathbb Z)=
\begin{cases}
\mathbb Z,&k=0,8,\\
\mathbb Z^4,&k=4,\\
0,&k=1,2,3,5,6,7.
\end{cases}
\]
In particular it is three-connected.
\end{theorem}
\begin{proof}
\begin{enumerate}[label=\arabic*.]
\item On the punctured base the evaluation map is
\[
 E_c:\Lambda_{\mathbb R}\longrightarrow F_c^\vee,
 \qquad E_c(x)(\omega)=\omega(x).
\]
Its kernel is zero: a real vector annihilating $F_c$ also annihilates
$\overline{F_c}$, and these two spaces span the complexified cohomology
space. Both real dimensions are six. If $\Omega$ is the evaluation
matrix of a local frame, its inverse is explicitly
\[
 E_c^{-1}(z)=R^{-1}
       \begin{pmatrix}\operatorname{Re}z\\\operatorname{Im}z\end{pmatrix},
 \qquad
 R=\begin{pmatrix}\operatorname{Re}\Omega\\\operatorname{Im}\Omega\end{pmatrix}.
\]
Under a frame change $\Omega'=G\Omega$, both $R$ and the real coordinate
column of $z$ are multiplied by the real matrix of $G$. Thus this
inverse glues. Evaluation sends each element of $\Lambda$ to its
specified period section and induces the marked real torus bundle
isomorphism, preserving the zero section. For a peripheral transport
$\gamma$ with matrix $M$, the actual evaluation square is
\[
 E_{c,\gamma p}(Mx)=\Phi_\gamma E_{c,p}(x),
\]
where $\Phi_\gamma$ is the transport in the period quotient. Consequently
\[
 J_c=E_c^{-1}iE_c=R^{-1}J_{\rm std}R,
 \qquad
 J_{\rm std}=\begin{pmatrix}0&-I_3\\I_3&0\end{pmatrix},
 \qquad J_c^2=-I_6,
\]
and $J_c(\gamma p)=M J_c(p)M^{-1}$.

\item At a finite site put $u=r e^{2\pi ia}$ on the universal angular
cover and $L=\log r/(2\pi)$. The prescribed attaching translation in
Theorem~\ref{night:P1.10} becomes
\[
 E_c^{-1}\!\left(-\frac{\log u}{2\pi i}E_c(v)\right)
       =-av+LJ_c(u)v.
\]
The affine action is $(x,u)\mapsto(Mx+v/m,\zeta_m u)$, and $Mv=v$ by
Theorem~\ref{night:P1.1}. Write
\[
 \alpha_0(x,u)=x-av,
 \qquad h_c(u)=LJ_c(u)v,
 \qquad \alpha_c=\alpha_0+h_c.
\]
The angular shift is $a\mapsto a+1/m$, and the displayed transport of
$J_c$ gives the exact identities
\[
 h_c(\zeta_m u)=M h_c(u),\qquad
 \alpha_c(Mx+v/m,\zeta_m u)=M\alpha_c(x,u).
\]
Hence $h_c$ is a section of the real local-system vector bundle on
the attaching annulus. Choose a scalar cutoff $\chi$ which is one at
the inner boundary of an attaching collar and zero at its outer
boundary. Translation by $-\chi h_c$ on that collar has inverse
translation by $\chi h_c$; the translations by $-s\chi h_c$,
$0\leq s\leq1$, give an equivariant isotopy. At the inner boundary
it carries $\alpha_c$ to $\alpha_0$ and at the outer boundary it is
the identity. The two finite collars are disjoint.

\item Write $L_k=\bigwedge^k\Lambda$ and $W_{i,k}=\bigwedge^kM_i$.
The cellular model of the main torus bundle over a graph with two
oriented edges has
\[
G_k=L_k\oplus L_{k-1}\oplus L_{k-1},\qquad
d_G(f,x_1,x_2)=((W_{1,k-1}-I)x_1+(W_{2,k-1}-I)x_2,0,0).
\tag{1.12.1}
\]
The boundary mapping torus at $i$ has
$B_{i,k}=L_k\oplus L_{k-1}$ and
\[
d_{B_i}(f,x)=((W_{i,k-1}-I)x,0).
\tag{1.12.2}
\]
The finite boundary maps to $G$ by identity on $L_k$ and inclusion
of the suspension into its respective edge.  The cusp boundary,
corresponding to $(g_1g_2)^{-1}$, maps by
\[
j_{0,G}(f,x)=(f,-W_{2,k-1}W_{0,k-1}x,-W_{0,k-1}x).
\tag{1.12.3}
\]
This follows by traversing the second edge backwards and then the
first backwards, transporting the fibre chain to each initial end.
Indeed its boundary is $W_{0,k-1}-I$ because
\[
-(W_1-I)W_2W_0-(W_2-I)W_0=W_0-I.
\]

\item Here are full models for the finite pieces and their attaching maps.
Put $\varepsilon_i=\ell(v_i)$ and
$P_i=(\varepsilon_iv_i,e_1,\ldots,e_5)$.
This is unimodular.  It splits
$\Lambda=\mathbb Z(\varepsilon_iv_i)\oplus K$, $K=\ker\ell$,
and $P_i^{-1}M_iP_i=1\oplus R_i$.
The cap retracts onto the mapping torus of $R_i$ on $T_K$.
Its chains are
\[
E_{i,k}=\bigwedge^kK\oplus\bigwedge^{k-1}K,
\qquad d_{E_i}(f,x)=((\bigwedge^{k-1}R_i-I)x,0).
\tag{1.12.4}
\]
In these split coordinates, with $v=\varepsilon_iv_i$, the finite
boundary-to-cap chain map is specified on all summands by
\[
\begin{array}{c|c}
\text{boundary chain}&\text{cap chain}\\\hline
x\in\bigwedge^kK&(x,0)\\
v\wedge x&(0,\varepsilon_i\sum_{j=0}^{m_i-1}
                  (\bigwedge^{k-1}R_i)^jx)\\
g_i\wedge x&(0,x)\\
g_i\wedge v\wedge x&0.
\end{array}
\tag{1.12.5}
\]
Changing back uses
$\operatorname{diag}(\bigwedge^kP_i^{-1},
\bigwedge^{k-1}P_i^{-1})$.
Geometrically this map is
$(t,y,\theta)\mapsto[y,\theta+\varepsilon_im_it]$ into the mapping
torus.  The $t$-circle sweeps $m_i$ strips, giving the norm sum; the
$g_i$-circle sweeps one strip; their product maps to a single
base direction, giving zero on the last row.  This proves the map
including its sign and the possible finite-cap torsion.

\item We justify that these are integral cellular models, rather than a
formal use of invariant rational forms.  In the bases with columns
\[
K_1=\begin{pmatrix}
-1&-1&0&-1&0\\0&1&0&1&1\\0&1&0&0&0\\
0&-1&-1&0&-1\\0&0&1&0&0
\end{pmatrix},\qquad
K_2=\begin{pmatrix}
-1&0&-1&0&-1\\-1&1&0&1&-1\\0&1&-1&1&-1\\
0&-1&0&-1&1\\1&-1&1&0&1
\end{pmatrix},
\]
both unimodular, $R_1$ is the permutation of five coordinates in a
cycle, and $R_2$ is a four-cycle together with multiplication by $-1$
on the fifth coordinate.  The minimal product cell structure of
circles is preserved by these signed permutations; its cellular
boundary is zero and its action on cells is exactly the exterior
power.  Formula (1.12.5) is obtained by subdividing its swept
base interval into the indicated strips.  The common torus can be
marked into these two structures by based cellular approximations.
Their cellular maps are the exterior-power matrices because the
cellular differentials vanish.  Based maps between these aspherical
tori with a fixed map on $\pi_1$ have a contractible space of
homotopies: lifting to the universal covers gives equivariant maps
to a contractible vector space and linear interpolation relative to
the base point.  Thus the choices introduce no extra attaching class
on the graph.  Composing the two edge cylinders gives exactly
(1.12.3).  The displayed conjugators are checked by exact multiplication in
\path{code/current_audits/audit_corrections/code/independent_V_topology_inputs_07.py}.
Its JSON output records the actual split bases, the two signed
permutation matrices and their cycles: one positive five-cycle, and
one positive four-cycle together with a negative fixed coordinate.

\item Let $F_k:B_{1,k}\oplus B_{2,k}\to G_k\oplus E_{1,k}\oplus E_{2,k}$
have the main boundary maps as its first row, and the negatives of
(1.12.5) in the corresponding two cap rows.  The chain complex of
$Y$ is the mapping cone
\[
C_k(Y)=G_k\oplus E_{1,k}\oplus E_{2,k}
       \oplus B_{1,k-1}\oplus B_{2,k-1},\qquad
d_Y=\begin{pmatrix}d_G\oplus d_{E_1}\oplus d_{E_2}&F_{k-1}\\
0&-d_{B_1}\oplus-d_{B_2}\end{pmatrix}.
\tag{1.12.6}
\]
Equations (1.12.1)--(1.12.6), minors for exterior powers, and the input
matrices specify every entry of every differential.
The full arrays are reconstructed and exported in
\path{code/current_audits/audit_corrections/code/reconstruct_V_global_06.json}
by its accompanying Python program.  It verifies $d^2=0$ in all
degrees and carries out integral Smith reductions.  The separate
\path{code/current_audits/audit_corrections/code/verify_V_global_saved_07.py}
reconstructs the homology bases used below and checks the saved
integral maps.  Chain ranks in degrees zero through eight are
$(3,22,71,132,155,118,57,16,2)$.
The resulting homology of $Y$ is free, with ranks
$(1,1,3,9,12,9,3,0,0)$.

\item The boundary $B_0$ has
\[
H_1(B_0)=\mathbb Z^4,\qquad H_2(B_0)=\mathbb Z^{12},\qquad
H_3(B_0)=\mathbb Z^{18}\oplus\mathbb Z/2.
\]
The homomorphism $j_Y:H_k(B_0)\to H_k(Y)$ from (1.12.3) is
surjective for $k=1,2,3$.  In degrees two and three, the saved-output verifier supplies integral
cycle bases of its free kernel represented in $L_k$ itself.  It
exports their full fibre coordinates and checks both their classes
in $H_k(B_0)$ and their images in $H_k(Y)$.  Their entire suspension
block in (1.12.2) is exactly zero over $\mathbb Z$.  In degree three
the additional order-two generator is also given as a fibre cycle
and its cusp image is checked to be zero.  Therefore these kernels are represented by cycles of
the nearby torus itself.  The restriction of the actual map to the
cusp cap on these kernels is consequently the already proved
collapse of P1.11, regardless of its value on other boundary
suspension cycles.

\item Use the homology bases and actual cycle representatives exported in
\path{code/current_audits/audit_corrections/code/verify_V_global_saved_07.json}.
The boundary bases are computed from (1.12.2), the complement bases
from (1.12.6), and the cusp bases are those of the marked collapse
in P1.11.  In these bases the combined map in degree one is
\[
\begin{pmatrix}20&1&0&0\\1&0&-1&-1\\1&0&0&0\\0&0&2&1\end{pmatrix},
\qquad\det=1.
\tag{1.12.7}
\]
The cusp meridian maps to zero here because the zero section extends.
In degree two $\ker j_Y=\mathbb Z^9$ and its map to $H_2(C)$ is
\[
\begin{pmatrix}
-1&0&0&-1&-1&0&0&0&0\\
0&0&1&1&0&0&-1&1&1\\
0&0&0&0&1&0&0&2&1\\
0&0&0&0&0&0&0&-1&-1\\
0&0&0&0&1&0&0&0&0\\
0&0&0&-1&0&0&0&0&0\\
0&0&0&-1&0&-1&1&0&0\\
0&1&0&-1&0&0&1&0&0\\
0&1&0&-2&0&-1&1&0&0
\end{pmatrix},\qquad\det=1.
\tag{1.12.8}
\]
In degree three $\ker j_Y=\mathbb Z/2\oplus\mathbb Z^9$.
Its collapse matrix, with the torsion generator first, is
\[
\begin{pmatrix}
0&0&0&0&0&-1&0&-1&-1&1\\
0&0&0&0&0&0&0&1&1&0\\
0&0&1&-1&-1&0&-1&0&0&0\\
0&0&0&0&1&0&0&0&0&0\\
0&0&0&0&0&0&0&1&0&0\\
0&0&0&1&1&0&0&0&0&0\\
0&0&1&-2&-1&1&0&1&2&0\\
0&0&1&-2&-1&0&-1&0&0&1\\
0&1&0&0&0&0&1&0&0&1
\end{pmatrix}.
\tag{1.12.9}
\]
The determinant of its last nine columns is one.  In particular the
torsion of the boundary and the second finite cap has been included
throughout the calculation.

\item We now apply Mayer--Vietoris to $M_c=Y\cup C_{\rm cap}$.
Because $j_Y$ is onto, subtracting a preimage of the $Y$-coordinate
identifies the cokernel of $(j_Y,j_C)$ with
$H_k(C)/j_C(\ker j_Y)$; its kernel is
$\ker j_Y\cap\ker j_C$.
Equations (1.12.7)--(1.12.9) therefore say that the combined map is an
isomorphism in degrees one and two and is onto in degree three.
The exact sequence yields $H_2(M_c)=H_3(M_c)=0$.
We already proved $\pi_1(M_c)=0$, and successive applications of
Hurewicz now give $\pi_2(M_c)=\pi_3(M_c)=0$.

\item Integral Poincar\'e duality and the universal coefficient theorem
give $H_5=H_6=H_7=0$.  They also show that $H_4$ is free: its possible
torsion would be $\operatorname{Ext}(H_4,\mathbb Z)$ in $H^5$, which
vanishes by duality with $H_3=0$.
Finally $\chi=6$ gives $\operatorname{rank}H_4=4$.
This proves all claimed integral homology groups without making an
assumption about the higher-degree cusp suspension chain term.

\end{enumerate}
\end{proof}

% CITED INPUT, WITH ITS PROOF: night:P1.2
\begin{theorem}\label{night:P1.2}
\StatementRecord{night:P1.2}{proved}{}{Definitions and the displayed constructions}
The lower blocks of $A,B$ are precisely the quotient representations of
$\bigwedge^2(C_f^{-1}),\bigwedge^2C_g$ in the supplied basis $Q$.
The integral quotient is the dual of the primitive exterior-square
lattice.  The natural map from the primitive lattice into that quotient
has cokernel $\mathbb Z/2$; these two integral representations are not
isomorphic.  Their rationalizations are isomorphic by that same map.
\end{theorem}
\begin{proof}
\begin{enumerate}[label=\arabic*.]
\item The two matrices obtained from their two-by-two minors are
\[
Q^{-1}\bigwedge^2(C_f^{-1})Q=
\begin{pmatrix}
1&-1&0&1&0&0\\0&1&-1&-1&1&0\\0&1&0&-1&0&0\\
0&2&-1&-1&0&1\\0&2&0&-1&0&0\\0&1&1&-1&0&0
\end{pmatrix},
\]
\[
Q^{-1}\bigwedge^2C_gQ=
\begin{pmatrix}
1&0&1&0&0&0\\0&0&-1&0&1&0\\0&0&-1&0&0&1\\
0&0&-1&0&0&0\\0&1&-1&0&0&0\\0&0&-1&1&0&0
\end{pmatrix}.
\]
Their lower five-by-five blocks are the asserted blocks, entry by entry.
Put $L=\bigwedge^2\mathbb Z^4$, let $(\ ,\ )$ be its unimodular wedge
pairing, and set $\omega=(1,1,1,1,1,1)^t$.
Then $(\omega,\omega)=2$ and the primitive sublattice
$P_{\mathbb Z}=\omega^\perp$ is the kernel of
$(1,-1,1,1,-1,1)$.  In the basis with columns
\[
\begin{pmatrix}
1&-1&-1&1&-1\\1&0&0&0&0\\0&1&0&0&0\\
0&0&1&0&0\\0&0&0&1&0\\0&0&0&0&1
\end{pmatrix},
\]
the inclusion followed by projection
$j:P_{\mathbb Z}\longrightarrow L/\mathbb Z\omega$ has matrix
\[
j=\begin{pmatrix}
0&1&1&-1&1\\-1&2&1&-1&1\\-1&1&2&-1&1\\
-1&1&1&0&1\\-1&1&1&-1&2
\end{pmatrix},\qquad \det j=2.
\]
Its Smith factors are $(1,1,1,1,2)$.  The unimodular pairing on $L$
identifies $L/\mathbb Z\omega$ with $P_{\mathbb Z}^\vee$;
under this identification $j$ is the pairing map.

\item We rule out every possible integral intertwiner.  Denote the lower blocks above by $A_5,B_5$.
The Krylov matrix $(e_3,A_5e_3,A_5^2e_3,A_5^3e_3,A_5^4e_3)$
has determinant $-1$, so $A_5$ is cyclic.  Any rational endomorphism
commuting with $A_5$ therefore has the form
$u_0I+u_1A_5+\cdots+u_4A_5^4$.
The entries $(1,2),(1,4),(2,4),(3,1)$ of its commutator with $B_5$
give, on $(u_1,u_2,u_3,u_4)^t$, the coefficient matrix
\[
\begin{pmatrix}
0&1&0&1\\1&-1&0&0\\1&-1&-1&1\\1&1&0&1
\end{pmatrix},\qquad \det=-1.
\]
Thus $u_1=u_2=u_3=u_4=0$, proving that the simultaneous rational
centralizer is exactly $\mathbb QI$.
The natural map $j$ above is equivariant and rationally invertible.
For any rational equivariant map $f:P_{\mathbb Q}\to
(L/\mathbb Z\omega)_{\mathbb Q}$, the composite $fj^{-1}$ lies
in this centralizer.  Consequently $f=cj$ for some $c\in\mathbb Q$.
An integral isomorphism would have determinant $\pm1$, whereas
$\det(cj)=2c^5$.  Its $2$-adic valuation would give
$1+5v_2(c)=0$, impossible because $v_2(c)\in\mathbb Z$.
This proves the asserted failure of integral isomorphism and retains
$j$ itself as the explicit rational isomorphism.

\end{enumerate}
\end{proof}

% CITED INPUT, WITH ITS PROOF: night:P1.7
\begin{theorem}\label{night:P1.7}
\StatementRecord{night:P1.7}{proved}{}{\cref{maps:src:F.3.5}, \cref{night:P1.2}}
The integral $H^1$ monodromy of the pencil has a companion marking with
finite pair $(C_f^{-1},C_g)$, and cusp product $C_fC_g^{-1}$.
Its exterior-square quotient is consequently the integral quotient
representation appearing in P1.2.
\end{theorem}
\begin{proof}
\begin{enumerate}[label=\arabic*.]
\item Let $L=H^1(C_\lambda,\mathbb Z)$ with its integral cup pairing,
and write $T_1,T_2$ for the finite generators of
\cref{maps:src:F.3.5}. Put $A=C_f$ and $B=C_g$.
The rational intertwiner proved there is an actual map
\[
 q:L_{\mathbb Q}\longrightarrow\mathbb Q^4,
 \qquad qT_1q^{-1}=A^{-1},\quad qT_2q^{-1}=B.
\]
Its relation to the Levelt presentation is
$(h_\infty,h_1,h_0)=(T_1^{-1},T_1T_2,T_2^{-1})$;
the ordered product is one. The known rigidity input and its explicit
rational change of basis are recorded in \cref{maps:src:F.3.5}.
The nodal peripheral is $c=(T_1T_2)^{-1}$. Its primitive
Picard--Lefschetz class $\delta\in L$ and coefficient one are the
geometric node input proved there: the nonseparating vanishing circle
has a transverse integral cycle meeting it once, and the annulus
$xy=t$ turns once. Thus $T_1T_2$ acts by
$v\mapsto v+\epsilon\delta\langle\delta,v\rangle$,
where $\epsilon\in\{1,-1\}$ records the fixed inverse and pairing
conventions.

\item A $B$-invariant alternating form has upper entries
$(u,v,u,u,v,u)$. Invariance under $A$, evaluated on the first and
fourth basis vectors, gives $v=-u$. The common invariant forms are
therefore exactly the multiples of
\[
 J=\begin{pmatrix}
 0&1&-1&1\\-1&0&1&-1\\1&-1&0&1\\-1&1&-1&0
 \end{pmatrix},\qquad
 d=\begin{pmatrix}1\\1\\1\\0\end{pmatrix}.
\]
The Pfaffian of $J$ is $1$, and direct multiplication gives
\[
 A^tJA=B^tJB=J,\qquad
 A^{-1}B-I=d e_4^t,\qquad d^tJ=e_4^t.
\]
Here $e_4$ is the fourth standard vector. Write the transported cup
form as $sJ$ and $q(\delta)=k d$, with $s,k\in\mathbb Q^*$.
Comparison with the actual Picard--Lefschetz formula gives
$\epsilon s k^2=1$.

\item The integral classes
$\delta_i=T_1^{-i}\delta$, $0\leq i\leq3$, have rational images
equal to the columns of $kD$, where
\[
 D=(d,Ad,A^2d,A^3d)=
 \begin{pmatrix}1&0&-1&0\\1&1&-1&-1\\1&1&0&-1\\0&1&0&0\end{pmatrix},
 \qquad \det D=1.
\]
Their actual pairing is $\epsilon G$, with
\[
 G=D^tJD=
 \begin{pmatrix}0&1&0&0\\-1&0&1&0\\0&-1&0&1\\0&0&-1&0\end{pmatrix},
 \qquad \det G=1.
\]
These four classes are a basis of $L_{\mathbb Q}$. For any $v\in L$,
write $v=\sum_i x_i\delta_i$ with $x_i\in\mathbb Q$.
The vector $(\langle\delta_i,v\rangle)_i=\epsilon Gx$ is integral.
Since $G^{-1}$ is integral, so is $x$. Thus the four classes are an
integral basis of $L$.

\item The required integral marking is the precise map
$\mu=k^{-1}q:L\longrightarrow\mathbb Z^4$.
Indeed $q(L)=kD\mathbb Z^4=k\mathbb Z^4$. Its inverse is specified
on standard vectors by
\[
 \mu^{-1}(e_j)=\sum_{i=0}^3(D^{-1})_{i+1,j}\delta_i,
 \qquad
 D^{-1}=\begin{pmatrix}1&-1&1&0\\0&0&0&1\\0&-1&1&0\\1&-1&0&1\end{pmatrix}.
\]
This supplies both integral arrows and retains the cup form
$\epsilon J$. In this marking the finite pair is $(A^{-1},B)$.

\item The exact simultaneous companion reversal is
\[
 W=\begin{pmatrix}0&0&0&1\\0&0&1&0\\0&1&0&0\\1&0&0&0\end{pmatrix}.
\]
It satisfies $W^2=I$, $\det W=1$, $WAW=A^{-1}$,
$WBW=B^{-1}$ and $W^tJW=-J$. These identities give the actual
integral comparison of the two companion conventions, including
the pairing sign. For the based cusp used in the statement take
$c'=T_1^{-1}cT_1=(T_2T_1)^{-1}$. Then
\[
 \mu c\mu^{-1}=B^{-1}A,\qquad
 \mu c'\mu^{-1}=A(B^{-1}A)A^{-1}=AB^{-1}.
\]
The homology matrices are the inverse transposes under the actual
evaluation pairing with $L$.

\item Apply $\bigwedge^2\mu$ to the integral wedge lattice. In the
lexicographic basis of \cref{night:P1.2}, the invariant vector is
$\omega=(1,1,1,1,1,1)^t$. Let $Q$ have first column $\omega$ and
remaining columns the last five standard vectors, as in that result.
The induced quotient is the actual map
\[
 [u\wedge v]\longmapsto
 \operatorname{pr}_{2,\ldots,6}
 Q^{-1}\bigl(\mu(u)\wedge\mu(v)\bigr).
\]
Its kernel before passage to the quotient is precisely the invariant
line. Its inverse is induced by $(\bigwedge^2\mu^{-1})Q$ on the last
five standard vectors. The finite quotient matrices are the lower
five-by-five blocks of
\[
 Q^{-1}(\bigwedge^2 A^{-1})Q,\qquad
 Q^{-1}(\bigwedge^2 B)Q.
\]
They are exactly the integral quotient matrices of
\cref{night:P1.2}, and the cusp follows by the same based word
$(T_2T_1)^{-1}$. The invariant form, actual Picard--Lefschetz rank-one map,
unimodular orbit matrix, integral inverse, reversal, based peripheral
word and exterior-quotient identities are recomputed exactly in
\path{code/current_audits/audit_corrections/code/independent_V_markings_06.py}; its JSON output records those checks.
The integral saturation argument in steps 2--4 supplies the map $\mu$
for the actual geometric lattice.
\end{enumerate}
\end{proof}
